\documentclass[12pt, letterpaper]{article}
\usepackage{graphicx} % Required for inserting images
%\usepackage[T1]{fontenc}

\usepackage{amsfonts}
\usepackage{tabularx}
\usepackage{multirow}
\usepackage{booktabs}
\usepackage{float}
%\usepackage{natbib}

\usepackage{polski}
%\usepackage[utf8]{inputenc}

\title{Programy użytkowe \\ Lista 4}

\date{}

\begin{document}
\maketitle
\begin{enumerate}
    \item Użyj pakietów \verb|hyperref| oraz polecenia \verb|\tableofcontents| do utworzenia spisu treści na początku dokumentu, który zawiera co najmniej dwie sekcje oraz po jednej podsekcji.
    \item Dodaj w polu autora afiliację, używając znaku nowej linii lub korzystając z pakietu \verb|authblk|.
    \item Dodaj streszczenie notatki (\verb|abstract|).
    \item Używając \verb|marginpar| nanieś komentarz na marginesie akapitu oraz równania w trybie wyróżnionym.\marginpar{Efekt ma być taki.}
    \item Używając polecenia \verb|\footnote| dodaj komentarz w stopce.
    \item Utwórz bibliografię na końcu notatki. Niech zawiera ona co najmniej dwa elementy, jeden będący artykułem, a drugi książką, na przykład
   % \bibliographystyle{plainnat}
    \begin{center}
    \begin{minipage}{0.7\linewidth}
        \begin{thebibliography}{9}
 %   \bibliographystyle{plainnat}
    \bibitem{Einstein:1905ve}
A.~Einstein,
\textit{On the electrodynamics of moving bodies},
Annalen Phys. \textbf{17} (1905), 891-921
%681 citations counted in INSPIRE as of 27 Oct 2025
\bibitem{Dirac30}P. A. M. Dirac, \textit{The Principles of Quantum Mechanics}, Clarendon Press, 4th edition, [1930] 1981
\end{thebibliography}
    \end{minipage}
        \end{center}
        Użyj alternatywnej metody tworzenia odnośników, np. [Ein05] zamiast [1] i [Dir30] zamiast [2].
    \item Korzystając z polecenia \verb|\appendix| dołącz dwie sekcje uzupełnienia na końcu notatki, za bbliografią. 
\end{enumerate}
\textbf{Zadanie domowe.} Zapisz w \LaTeX~zadanie rozwiązanie przez Ciebie przy tablicy, albo fragment wykładu (twierdzenie, lemat, przykład).


\end{document}